\documentclass[12pt, a4paper]{article}

% 引入必要的宏包
\usepackage[UTF8]{ctex}     % 中文支持
\usepackage{amsmath, amssymb, amsthm} % 数学符号和环境
\usepackage{geometry}       % 页面设置
\usepackage{enumitem}       % 列表定制
\usepackage{fancyhdr}       % 页眉页脚
\usepackage{cases}          % 分段函数增强

% 页面边距设置
\geometry{left=2.5cm, right=2.5cm, top=2.5cm, bottom=2.5cm}

% 宏定义：方便排版选择题选项
% 横向排列的选项 (4列)
\newcommand{\choicesFour}[4]{
    \begin{enumerate}[label=\Alph*., itemsep=0pt, parsep=0pt, topsep=5pt, labelsep=0.5em]
        \begin{minipage}{0.22\linewidth} \item #1 \end{minipage}
        \begin{minipage}{0.22\linewidth} \item #2 \end{minipage}
        \begin{minipage}{0.22\linewidth} \item #3 \end{minipage}
        \begin{minipage}{0.22\linewidth} \item #4 \end{minipage}
    \end{enumerate}
}
% 横向排列的选项 (2列)
\newcommand{\choicesTwo}[4]{
    \begin{enumerate}[label=\Alph*., itemsep=0pt, parsep=0pt, topsep=5pt, labelsep=0.5em]
        \begin{minipage}{0.45\linewidth} \item #1 \end{minipage}
        \begin{minipage}{0.45\linewidth} \item #2 \end{minipage}
        \begin{minipage}{0.45\linewidth} \item #3 \end{minipage}
        \begin{minipage}{0.45\linewidth} \item #4 \end{minipage}
    \end{enumerate}
}
% 纵向排列的选项 (1列)
\newcommand{\choices}[4]{
    \begin{enumerate}[label=\Alph*., itemsep=0pt, parsep=0pt, topsep=5pt, labelsep=0.5em]
        \item #1
        \item #2
        \item #3
        \item #4
    \end{enumerate}
}

% 填空题下划线
\newcommand{\blank}[1]{\underline{\hspace{#1}}}

\title{\textbf{《高等数学》必刷习题——提高版}}
\author{}
\date{}

\begin{document}

\section*{第十章 向量代数与空间解析几何（提高）}

\begin{enumerate}
    \item 直线 $ L_{1} $ : $ x - 1 = \frac{y - 5}{-2} = z + 8 $ 与直线 $ L_{2} $ : $ \begin{cases} x - y = 6 \\ 2y + z = 3 \end{cases} $ 的夹角为 \blank{1cm}。
    \choicesFour{$ \frac{\pi}{6} $}{$ \frac{\pi}{4} $}{$ \frac{\pi}{3} $}{$ \frac{\pi}{2} $}

    \item 曲面 $ z = 4 - x^{2} - y^{2} $ 上一点 $P$ 的切平面平行于平面 $ 2x + 2y + z - 1 = 0 $ ，则 $P$ 点坐标为 \blank{1cm}。
    \choicesFour{$ (1,-1,2) $}{$ (-1,1,2) $}{$ (1,1,2) $}{$ (-1,-1,2) $}

    \item 直线 $l$: $ \frac{x+3}{-2}=\frac{y+4}{-7}=\frac{z}{3} $ 与平面 $ \pi $: $4x-2y-2z-3=0$ 的位置关系是 \blank{1cm}。
    \choicesFour{平行}{垂直相交}{$l$ 在 $ \pi $ 上}{相交但不垂直}

    \item 已知向量 $\mathbf{a}$ 垂直于向量 $ \mathbf{b} = 2\mathbf{i} - 3\mathbf{j} + \mathbf{k} $ 和 $ \mathbf{c} = \mathbf{i} - 2\mathbf{j} + 3\mathbf{k} $ ，且满足于 $ \mathbf{a} \cdot (\mathbf{i} + 2\mathbf{j} - 7\mathbf{k}) = 10 $ ，求 $\mathbf{a} =$ \blank{1cm}。
    \choicesTwo
    {$ -7\mathbf{i} - 5\mathbf{j} - \mathbf{k} $}
    {$ 7\mathbf{i} + 5\mathbf{j} + \mathbf{k} $}
    {$ -5\mathbf{i} - 3\mathbf{j} - \mathbf{k} $}
    {$ 5\mathbf{i} + 3\mathbf{j} + \mathbf{k} $}

    \item 向量 $\mathbf{a}$ 与 $x$ 轴与 $y$ 轴构成等角，与 $z$ 轴夹角是前者的 2 倍，下面哪一个代表的是 $\mathbf{a}$ 的方向 \blank{1cm}。
    \choicesTwo
    {$ \alpha=\frac{\pi}{2},\beta=\frac{\pi}{2},\gamma=\frac{\pi}{4} $}
    {$ \alpha=\frac{\pi}{4},\beta=\frac{\pi}{4},\gamma=\frac{\pi}{8} $}
    {$ \alpha=\frac{\pi}{4},\beta=\frac{\pi}{4},\gamma=\frac{\pi}{2} $}
    {$ \alpha=\pi,\beta=\frac{\pi}{2},\gamma=\frac{\pi}{2} $}

    \item 设已知两点 $A(4,0,5)$ 和 $B(7,1,3)$，方向和 $ \overrightarrow{AB} $ 一致的单位向量为 \blank{2cm}。

    \item 已知 $ \mathbf{a} = \{2, 1, 2\} $ , $ \mathbf{b} = \{4, -1, 10\} $ , $ \mathbf{c} = \mathbf{b} - \lambda \mathbf{a} $ , 且 $ \mathbf{a} \perp \mathbf{c} $ , 则 $ \lambda = $ \blank{2cm}。

    \item 过点 $ (3,2,1) $ 与直线 $ x = y = z $ 垂直的平面方程为 \blank{2cm}。

    \item 过点 $ (3,-2,2) $ 与直线 $ x=\frac{y}{2}=z $ 垂直的平面方程是 \blank{2cm}。

    \item 过点 $ (2,2,3) $ 且与平面 $ 3x + y = z + 2 $ 垂直的直线方程是 \blank{2cm}。

    \item 直线 $ \frac{x-1}{2} = \frac{y+2}{5} = \frac{z-4}{-3} $ 的方向向量 $s=$ \blank{1.5cm}，与平面 $ 2x + 5y - 3z - 4 = 0 $ 是 \blank{1.5cm} 的。

    \item 求平行于 $y$ 轴且过点 $ P(1,5,1) $ , $ Q(3,2,-1) $ 的平面方程。

    \item 已知 $ \triangle ABC $ 的三个顶点的坐标分别为 $ A(1,2,3) $ , $ B(3,4,5) $ , $ C(2,4,7) $ ,
    求：(1) $ \triangle ABC $ 所在的平面方程；
    (2) $AB$ 边上的中线 $CD$ 所在的直线方程。

    \item 求过点 $ (-1,-4,3) $ 并与两直线 $ L_{1}\begin{cases}2x-4y+z=1\\ x+3y+5=0\end{cases} $ 和 $ L_{2}\begin{cases}x=2+4t\\ y=-1-t\\ z=-3+2t\end{cases} $ 都垂直的直线方程。

    \item 求通过点 $ M_{1}(3,-5,1) $ ， $ M_{2}(4,1,2) $ 且垂直于平面 $ x-8y+3z-1=0 $ 的平面方程。
\end{enumerate}

\end{document}
